%% Q1
The features that really distinguish online testing from offline testing are:
\begin{itemize}
\item (1) offline testing is \emph{local} (conducted at the centre), while online testing is \emph{centralised} (server-based);
\item (4) the internet is a must for online testing but not for offline testing;
\item (5) online testing usually gives \emph{instant} evaluation, while offline testing need not.
\end{itemize}
Statement (2) does not hold, since offline tests may also be digital (e.g.\ OMR). Statement (3) is false, since both modes can be subjective or objective. The answer is 1, 4 and 5.

%% Q2
In Kilpatrick's \textbf{project method} (based on Dewey's pragmatism) each learner \emph{plans, executes and evaluates} a purposeful activity in a real-life setting. Its strength lies in the learner's \textbf{individualised, active engagement}: learning by doing. The teacher acts as a guide rather than an instructor.

%% Q3
Computer-based testing (CBT) makes testing \textbf{transparent, objective and efficient}: standardised administration, randomised question sets, automatic scoring, quick results and secure records. The other options are narrower by-products of this main purpose.

%% Q4
An \textbf{open-book examination} lets candidates consult books. It therefore does not test memory; it tests the \textbf{ability to apply knowledge}: analysis, interpretation and problem solving.

%% Q5
\textbf{Reflective teaching} makes learners think about what they learn and how. They \textbf{internalise} knowledge and develop the ability to \emph{judge} and evaluate critically (the reflective level of learning in Hunt's model: memory $\rightarrow$ understanding $\rightarrow$ reflective).

%% Q6
In the \textbf{experimental method} the researcher manipulates independent variables and controls extraneous ones, often in a laboratory. The data therefore come from a \textbf{contrived, not a natural setting}. This is its main limitation (low ecological validity).
\why{Experimental data actually involve \emph{precise} measurement, so (B) is wrong. ``Unrealistic'' and ``artificial'' overstate the point; the defining feature is the controlled rather than natural setting.}

%% Q7
Historical research depends on records of the past. Its biggest constraint is the \textbf{non-availability of primary sources} (original documents, eyewitness accounts), which are often lost, destroyed or inaccessible. Researchers then have to rely on secondary sources.

%% Q8
The title and hypotheses are formulated \emph{before} data collection and guide the entire study. Changing them after the data have been analysed, to fit the results, amounts to HARKing (Hypothesising After the Results are Known). It is \textbf{illogical and not desirable} and undermines the scientific integrity of the research.

%% Q9
A research title is a precise statement of what the study is about: the variables, population and relationship being investigated. As per the official key, the title is treated as an \textbf{assertion}, a declarative statement of the research problem.

%% Q10
The logical sequence of scientific research:
\begin{center}
\begin{tikzpicture}[node distance=3mm,every node/.style={draw=brand,rounded corners,fill=soft,font=\scriptsize,align=center,text width=15mm,minimum height=9mm}]
\node (a) {(3) Identify problem}; \node[right=of a] (b) {(1) Formulate hypothesis}; \node[right=of b] (c) {(2) Develop tool}; \node[right=of c] (d) {(4) Select sample}; \node[right=of d] (e) {(5) Analyse \& interpret};
\foreach \x/\y in {a/b,b/c,c/d,d/e} \draw[-Stealth,brand] (\x)--(\y);
\end{tikzpicture}
\end{center}
The order is \textbf{3, 1, 2, 4, 5}.

%% Q11 ctx
\textbf{Directions for Q.~11 to 15:} Read the following passage and answer the questions.\par\smallskip
If the world is different today, we need to think, talk and engage accordingly. Falling back is unlikely to help, he said adding ``purposeful pursuit of national interest is shifting global dynamics''. Highlighting India's new approach in dealing with terrorism he compared ``lack of response to the terror attack of Mumbai in comparison to the way the country responded to the Uri and Pulwama strikes. On India walking away from the Regional Comprehensive Economic Partnership, the external affair minister said no agreement was better than having a bad agreement.

Giving a historical perspective to geopolitical issues, Jai-Shankar said ``For years India's position on the world state seemed assured but the 1962 conflict with China significantly damaged India's standing.

%% Q11
\textbf{RCEP} is the \textbf{Regional Comprehensive Economic Partnership}, a free-trade agreement among the 10 ASEAN countries and China, Japan, South Korea, Australia and New Zealand (signed 2020). India withdrew from the negotiations in November 2019.

%% Q12
The minister said that ``\textbf{no agreement was better than having a bad agreement}''. Option (C) expresses this: there is nothing wrong in saying no to an unfavourable agreement rather than accepting an awful one.

%% Q13
The passage names \textbf{(S.) Jaishankar}, India's External Affairs Minister, as the speaker. Option (B) (``Subhramanyam Jayshankar'') refers to Dr.~Subrahmanyam Jaishankar.

%% Q14
The passage contrasts the ``lack of response'' to the 2008 Mumbai attack with India's response to the Uri (2016, surgical strikes) and Pulwama (2019, Balakot air strike) attacks. The approach has changed from \textbf{reactive to proactive}.

%% Q15
The minister was ``giving a historical perspective to \textbf{geopolitical} issues'' when he spoke of the 1962 conflict with China, so the conflict is described as a geopolitical issue.

%% Q16
In the communication process the \textbf{sender encodes} the message into symbols suited to the medium, and the \textbf{receiver decodes} it to extract the meaning. The medium only carries the message.

%% Q17
\textbf{Proxemics} deals with the use of \textbf{space} and distance in communication.
\why{Appearance belongs to artifactual communication; touch is haptics; movements of body organs are kinesics.}

%% Q18
Every person has a larger \emph{recognition} vocabulary (words understood) than \emph{use} vocabulary (words actually used). A language barrier arises when the sender's \textbf{use vocabulary} contains words that are not in the receiver's \textbf{recognition vocabulary}, i.e.\ the two differ.

%% Q19
In \textbf{intrapersonal communication} a person communicates \emph{with oneself} (thinking, self-talk, reflection), so the same person is both sender and receiver.

%% Q20
Mass media shape public opinion, modify behaviour and change values and beliefs. \textbf{Judicial decisions} must be based on law and evidence and should not be influenced by the media, so this is the exception.

%% Q21
The pattern is $\times2+1, \times2+2, \times2+3, \ldots$:
\[4\xrightarrow{\times2+1}9\xrightarrow{\times2+2}20\xrightarrow{\times2+3}43\xrightarrow{\times2+4}\mathbf{90}\]
90 is not among the options, so the agency \textbf{dropped} the question.

%% Q22
$2P+3Q=77$ and $3P+5Q=124$.
Multiply the first equation by 5 and the second by 3: $10P+15Q=385$ and $9P+15Q=372$. Subtracting gives $P=\mathbf{13}$ (and $Q=17$).
Check: $2(13)+3(17)=77$ \checkmark; $3(13)+5(17)=124$ \checkmark.

%% Q23
Sugar $=40\%\times300=120$ g. Let $x$ g of sugar be added.
\[\frac{120+x}{300+x}=\frac12\Rightarrow240+2x=300+x\Rightarrow x=\mathbf{60\ g}\]

%% Q24 fig
\pyr{29}{16}{13}{27}{9}{33}\qquad\pyr{38}{21}{17}{35}{16}{43}\qquad\pyr{55}{30}{25}{?}{?}{?}

%% Q24
Two rules hold in every box:
\begin{itemize}
\item \textbf{Bottom middle} $=$ (difference of the middle row)$^2$: $(16-13)^2=9$, $(21-17)^2=16$, so $(30-25)^2=\mathbf{25}$.
\item \textbf{Bottom left $+$ bottom right} $=2\times\text{top}+2$: $27+33=60=2(29)+2$, $35+43=78=2(38)+2$, so the sum must be $2(55)+2=112$. Also bottom right $-$ bottom left $=6, 8, \ldots$, so $\mathbf{10}$.
\end{itemize}
Left $+$ right $=112$ and right $-$ left $=10$ give \textbf{51 and 61}. The missing numbers are \textbf{51, 25, 61}.

%% Q25
CP $=\dfrac{1536}{0.80}=1920$. New SP $=2000$, so profit $=80$.
\[\text{Profit \%}=\frac{80}{1920}\times100=\mathbf{4.17\%}\]

%% Q26
\textbf{Upamana} is knowledge by \textbf{comparison} or analogy. For example, a person told that ``a gavaya (wild ox) resembles a cow'' recognises a gavaya on seeing one.
\tip{Pratyaksha = perception; Anumana = inference; Upamana = comparison; Shabda = verbal testimony; Arthapatti = postulation; Anupalabdhi = non-perception.}

%% Q27
An odometer is the instrument that measures \emph{mileage} (distance travelled); a compass is the instrument that shows \textbf{direction}. The relation is \emph{instrument : what it measures}.

%% Q28 opts
\tikz[baseline=-0.5ex,scale=0.3]{\draw (0,0) circle(1.3); \draw (-0.55,0) circle(0.5); \draw (0.55,0) circle(0.5);}
\tikz[baseline=-0.5ex,scale=0.3]{\draw (0,0) circle(0.8); \draw (1.1,0) circle(0.8); \draw (2.2,0) circle(0.8);}
\tikz[baseline=-0.5ex,scale=0.3]{\draw (0,0) circle(1.1); \draw (0,0.1) circle(0.5); \draw (1.6,0) circle(0.9);}
\tikz[baseline=-0.5ex,scale=0.3]{\draw (0,0) circle(0.8); \draw (1.1,0) circle(0.8); \draw (3,0) circle(0.8);}

%% Q28
Lions and elephants are both \textbf{animals}, and no lion is an elephant. So the two small circles must be \textbf{separate} and both lie \textbf{inside} the large ``Animals'' circle, as in diagram (A).
\begin{center}
\begin{tikzpicture}[scale=0.7,font=\scriptsize]
\draw[fill=brand!8] (0,0) circle(2.1); \node at (0,1.6) {Animals};
\draw[fill=white] (-0.95,-0.2) circle(0.85) node {Lions}; \draw[fill=white] (0.95,-0.2) circle(0.85) node {Elephants};
\end{tikzpicture}
\end{center}

%% Q29
The Nyaya syllogism (\emph{Panchavayava}) has \textbf{five} members: Pratijna (proposition), Hetu (reason), Udaharana (example), Upanaya (application) and Nigamana (conclusion).

%% Q30
Both arguments are \textbf{strong}. Tall chimneys disperse pollutants high up and so reduce ground-level concentration, which directly protects people (I). But they spread the pollution into the upper atmosphere and over larger areas, contributing to acid rain and similar problems (II). Each argument is important and directly related to the question.

%% Q31
\[(a+b+c)^2=(a^2+b^2+c^2)+2(ab+bc+ca)=138+2(131)=400\ \Rightarrow\ a+b+c=\mathbf{20}\]

%% Q32 fig
\begin{tabular}{lc}\toprule \textbf{Wages (Rs.)} & \textbf{No.\ of workers}\\\midrule
Less than 50 & 4\\ Less than 60 & 10\\ Less than 70 & 20\\ Less than 80 & 40\\ Less than 90 & 50\\\bottomrule\end{tabular}

%% Q32 after
What percentage of workers earn wage more than Rs.~70?

%% Q32
This is a cumulative (``less than'') distribution. Total workers $=50$, and those earning less than Rs.~70 number 20.
Workers earning more than Rs.~70 $=50-20=30$, and $\dfrac{30}{50}\times100=\mathbf{60\%}$.

%% Q33 fig
\begin{tikzpicture}
\begin{axis}[ybar,width=11cm,height=4.6cm,bar width=12pt,ymin=0,ymax=28000,
  symbolic x coords={Gujarat,Haryana,Karnataka,Maharashtra,Rajasthan,Tamil Nadu,U.P.,Others},xtick=data,
  x tick label style={rotate=35,anchor=east,font=\scriptsize},yticklabel style={font=\scriptsize},
  nodes near coords,every node near coord/.append style={font=\tiny},
  title={\scriptsize Production of Rose in Various States},axis x line*=bottom,axis y line*=left,
  scaled y ticks=false]
\addplot[fill=white,draw=ink] coordinates {(Gujarat,15000) (Haryana,12500) (Karnataka,12450) (Maharashtra,20000) (Rajasthan,12400) (Tamil Nadu,22500) (U.P.,22400) (Others,25000)};
\end{axis}
\end{tikzpicture}

%% Q33 after
Which of the following state(s) contribute(s) less than 10 per cent in the total Rose production?

%% Q33
Total $=15000+12500+12450+20000+12400+22500+22400+25000=142{,}250$, so 10\% $=14{,}225$.
States below 14,225: \textbf{Haryana} (12,500), \textbf{Karnataka} (12,450) and \textbf{Rajasthan} (12,400). Gujarat (15,000) is above 10\%.

%% Q34 fig
\begin{minipage}{0.45\linewidth}\centering
{\scriptsize\bfseries Percentage of students (Total 4500)}\\[2pt]
\begin{tikzpicture}\pie[radius=1.35,text=pin,font=\tiny,color={brand!10,brand!30,brand!15,brand!40,brand!20,brand!5}]{11/Finance,22/HR,16/Mktg,18/IB,12/Banking,21/IT}\end{tikzpicture}
\end{minipage}\hfill
\begin{minipage}{0.45\linewidth}\centering
{\scriptsize\bfseries Ratio of Men to Women}\\[2pt]\small
\begin{tabular}{lcc}\toprule Specialization & Men & Women\\\midrule
HR & 2 & 3\\ Marketing & 7 & 5\\ IB & 5 & 1\\ Banking & 17 & 10\\ IT & 1 & 2\\ Finance & 5 & 4\\\bottomrule\end{tabular}
\end{minipage}

%% Q34 after
The number of women studying Finance is what percentage of the number of men studying the same?

%% Q34
Only the Finance ratio is needed: Men : Women $=5:4$.
\[\frac{\text{Women}}{\text{Men}}\times100=\frac45\times100=\mathbf{80\%}\]
(Numerically: Finance $=11\%$ of 4500 $=495$, of whom men $=275$ and women $=220$; $220/275=80\%$.)

%% Q35 fig
\begin{tabular}{cc}\toprule \textbf{Weight (kg)} & \textbf{No.\ of students}\\\midrule
35--38 & 3\\ 38--41 & 13\\ 41--44 & 19\\ 44--47 & 10\\ 47--50 & 5\\\bottomrule\end{tabular}

%% Q35 after
What is the average weight of the class?

%% Q35
Use class mid-points $x$:
\begin{center}\small
\begin{tabular}{cccc}\toprule Class & $x$ & $f$ & $fx$\\\midrule
35--38 & 36.5 & 3 & 109.5\\ 38--41 & 39.5 & 13 & 513.5\\ 41--44 & 42.5 & 19 & 807.5\\ 44--47 & 45.5 & 10 & 455.0\\ 47--50 & 48.5 & 5 & 242.5\\\midrule
& & 50 & 2128.0\\\bottomrule\end{tabular}
\end{center}
Mean $=2128/50=\mathbf{42.56}$ kg.

%% Q36
Internet Explorer (Microsoft, 1995--2022) is a \textbf{web browser}, software for accessing and viewing websites. Its successor is Microsoft Edge.

%% Q37
\textbf{Mozilla Firefox} is free and open-source (released under the Mozilla Public License), so its source code is publicly available to study and modify.
\why{Acrobat Reader is freeware but proprietary; MS PowerPoint and Windows are proprietary commercial software.}

%% Q38
\textbf{NMEICT} is the \textbf{National Mission on Education through Information and Communication Technology}, launched by MHRD in 2009. It supports e-content (e.g.\ NPTEL, e-PG Pathshala), connectivity and the SAKSHAT portal.

%% Q39
\textbf{ARPIT} (Annual Refresher Programme in Teaching), launched by MHRD in November 2018 on the SWAYAM platform, provides online refresher courses for higher-education faculty. Completing it counts for career advancement, like a conventional refresher course.
\why{MOODLE is an open-source LMS; NPTEL provides engineering and science courses; e-PG Pathshala provides PG e-content.}

%% Q40
LinkedIn, Instagram and Twitter (X) are social-networking platforms. \textbf{eBay} is an \emph{e-commerce} (online auction and shopping) website.

%% Q41
All four help abate air pollution: CNG vehicles emit fewer pollutants (i); avoiding firecrackers prevents spikes in particulate matter (ii); trees absorb pollutants and dust (iii); and not burning crop residue (stubble) prevents severe smog (iv). The answer is (i), (ii), (iii) and (iv).

%% Q42
The stratospheric ozone layer absorbs most of the Sun's harmful UV-B radiation. When it is depleted, \textbf{more ultraviolet radiation reaches the troposphere} and the Earth's surface, causing skin cancer, cataracts and crop damage.

%% Q43
NITI Aayog's \emph{Composite Water Management Index} report (2018) warned that 21 Indian cities, including Delhi, Bengaluru, \textbf{Chennai} and Hyderabad, could run out of groundwater by 2020. Among the options, Chennai is the South Indian city named. It faced an acute ``Day Zero'' water crisis in 2019.

%% Q44
India's AQI bands: Good (0--50), Satisfactory (51--100), \textbf{Moderate (101--200)}, Poor (201--300), Very poor (301--400), Severe (401--500). An AQI of 190 is \textbf{Moderate}.

%% Q45
\textbf{Poly lactic acid (PLA)} is made from renewable resources such as corn starch or sugarcane and is \textbf{biodegradable} and compostable, so it is environment-friendly.
\why{Polystyrene, polyethylene and polypropylene are petroleum-based, non-biodegradable plastics.}

%% Q46
\textbf{NAAC} (est.~1994 by UGC, Bengaluru) assesses and accredits higher-education institutions, so its purpose is \textbf{quality assurance}.
\why{Grants are given by UGC/HEGC, programmes are approved by statutory councils and universities, and universities are established by legislatures.}

%% Q47
The \textbf{AICTE} (All India Council for Technical Education) is the statutory body for technical education: planning, standards, approvals and so on.
\why{AIU is the Association of Indian Universities; NCTE regulates teacher education; AIIMS is a medical institute.}

%% Q48
\textbf{IMPRINT} (IMPacting Research INnovation and Technology), launched in 2015 by MHRD and led by the IITs and IISc, is a national initiative to address \textbf{major science and engineering challenges} in ten technology domains (health, energy, water, environment, security and others).

%% Q49
In \textbf{NIRF 2019} the \textbf{IIT Madras} ranked 1st in the Overall category, followed by IISc Bengaluru and IIT Delhi.

%% Q50
The \textbf{Kothari Commission} (Education Commission, 1964--66) recommended a uniform \textbf{10+2+3} pattern of education, which the National Policy on Education 1968 adopted.
\why{The Radhakrishnan Commission (1948--49) dealt with university education; the Balwantrai Mehta Committee (1957) dealt with Panchayati Raj; the Janardhan Reddy Committee (1992) reviewed the NPE 1986.}
